% =====================================================================
% 8. Practical conclusions and own activities
% =====================================================================
\section{Practical conclusions and my activities}\label{sec:conclusions}

\subsection{Practical conclusions}

\begin{enumerate}[leftmargin=1.4em,itemsep=1pt]
  \item \textbf{A shared, task-agnostic loop paid off.} The same five-step loop, logging and comparison
        tools served five classifier families, two generators and the transfer experiments; differences
        between blocks are therefore differences between models, and a bug fixed once was fixed
        everywhere.
  \item \textbf{Smoke runs and estimates before long runs.} Every long grid was preceded by a
        \texttt{--fast} run and, where it existed, an \texttt{estimate}. The two failures on the real
        hardware happened in steps not covered by them (generation memory, CLIP on the GPU): what is not
        tested on the target device is not tested.
  \item \textbf{Three seeds and paired intervals change the conclusions.} Several configurations that
        looked better by their mean were better in one seed of three only; without the paired bootstrap
        we would have reported them as improvements.
  \item \textbf{Training curves explain what the final metric hides:} faster or slower optimisation
        (batch norm and, weakly, LayerNorm placement; for skip connections the effect at our depth was too
        small to separate from other changes), overfitting (dropout), or a budget that was
        too short (the generators, the 10\,\% probe).
  \item \textbf{Generated numbers and checked quotes} made the reports slower to write and much easier to
        trust: a change in a log updates every report, and every literature claim has a page.
  \item \textbf{GPU memory is a design constraint.} The 4\,GB GPU fixed the batch of the transfer
        experiments, the context length of the generators and the generation batch; attention memory
        grows with the square of the sequence length.
\end{enumerate}

\subsection{My activities in the team}

Table~\ref{tab:contrib_of_mllb} gives my role in each completed task, generated from the task log
(\texttt{docs/PLAN\_DL.md}); ``did (shared)'' marks joint tasks.

\begin{table}[htbp]
\centering
\caption{My role in each completed task (generated from \texttt{docs/PLAN\_DL.md}).}
\label{tab:contrib_of_mllb}
% Generated by common/make_contributions.py from docs/PLAN_DL.md (OF). Do not edit.
{\small
\begin{tabularx}{\textwidth}{l>{\raggedright\arraybackslash}X>{\raggedright\arraybackslash}p{3.6cm}l}
\toprule
ID & Task & Role of Oriana Ferraz & Status \\
\midrule
T0 & Dataset shortlist (8 candidates) and choice & did (shared) & done \\
T1 & \texttt{DECISIONS.md} and \texttt{PLAN\_DL.md} first version & did (shared) & done \\
T2 & Decide the transfer-learning approach: own pre-trained model (not confirmed with the professor; raise it in the class discussion) & did (shared) & done \\
T3 & Verify licence and recommended folds on PhysioNet & did (shared) & done \\
T4 & Local project structure, \texttt{.gitignore}, environment file and GitHub repository & did & done \\
T5 & Download script, label construction, fold split, \texttt{Dataset}/\texttt{DataLoader} & reviewed & done \\
T6 & EDA: class balance, co-occurrence, demographics, signal quality, example traces & did & done \\
T7 & Evaluation module: macro AUROC, bootstrap CI, JSON logging, \texttt{make\_numbers.py} & reviewed & done \\
T8 & Shared training loop (zero\_grad, forward, loss, backward, step; train/eval modes) & did (shared) & done \\
T9 & Dense baseline + ablations & did & done \\
T10 & CNN 1D + ablations & did & done \\
T11 & LSTM/GRU (+bidirectional) and RNN + attention model + ablations & reviewed & done \\
T12 & Self-attention step by step (single layer, attention maps) and Transformer + ablations & reviewed & done \\
T13 & Transfer learning: ECG-GPT (T23) reused on the 5 superclasses, scratch vs frozen vs fine-tuning, 10\% and 100\% labels & did (shared) & done \\
T14 & Cross-cutting ablation: 12 leads vs single lead & did (shared) & done \\
T15 & Final comparison table and figures; comparison with published results & did (shared) & done \\
T16 & Literature: dataset and CNN-related papers & did & done \\
T17 & Literature: RNN/Transformer-related papers & reviewed & done \\
T18 & Learning-strategies annex of the PRDL reports (from checked quotes; the first class assignments were not available) & did her own part, cross review & done \\
T19 & PRDL-DL report (theory + critical analysis), opening with a summary of at most 3 pages & did her own part, cross review & done \\
T20 & MLLB-DL report (Section 0 AI environment, code, training, evaluation), opening with a summary of at most 3 pages & did her own part, cross review & done \\
T21 & \texttt{AI\_USAGE\_LOG.md} and AI Error Bounty sheets & did her own part, cross review & done \\
T22 & Adversarial review before delivery; preparation of the class discussion and oral defence & did (shared) & done \\
T23 & Generative autoregressive model (ECG-GPT): tokenisation of lead II, causal Transformer and LSTM, sampling with temperature, ablations (D13) & did & done \\
T24 & Zero-shot learning: theory and a short experiment on the class CLIP notebook & did (shared) & done \\
\bottomrule
\end{tabularx}}

\end{table}

On the engineering side I created the repository, the virtual environment and the folder layout, and I
ran every command of the project on my laptop. For the EDA, the dense network, the CNN and the
generators I specified the experiments, reviewed each module and its tests with the assistant before
the long runs, watched the smoke runs, fixed with the assistant what they revealed (Appendix~\ref{app:ai-code})
and checked the figures. The training loop and the experiments of area~2 were written and run together
with Mar\'ia, and I reviewed the code of her blocks; their implementation is described here from her
block notes and credited to her.
